% Part: incompleteness
% Chapter: incompleteness-provability
% Section: provability-conditions

\documentclass[../../../include/open-logic-section]{subfiles}

\begin{document}

\olfileid{inc}{inp}{prc}

\olsection{$\Th{PA}$-এর \usetoken{S}{derivability} শর্তগুলি}

পেয়ানো পাটিগণিত বা $\Th{PA}$ হলো সব !!{formula}s-এর জন্য আরোহ
স্বতঃসিদ্ধ যোগ করে পাওয়া $\Th{Q}$-এর প্রসারণ। অন্য কথায়, প্রতিটি
!!{formula}~$!A$-এর জন্য $\Th{Q}$-তে নিচের আকারের স্বতঃসিদ্ধ যোগ করা হয়:
\[
(!A(0) \land \lforall[x][(!A(x) \lif !A(x'))]) \lif \lforall[x][!A(x)]
\]
খেয়াল করো, এটি আসলে একটি \emph{ছক}; অর্থাৎ এতে অসীমসংখ্যক স্বতঃসিদ্ধ
আছে (এবং দেখা যায় $\Th{PA}$ সসীমভাবে স্বতঃসিদ্ধায়নযোগ্য {\em নয়})।
কিন্তু কোনো প্রতীকক্রম আরোহ স্বতঃসিদ্ধের দৃষ্টান্ত কি না কার্যকরভাবে
নির্ধারণ করা যায় বলে $\Th{PA}$-এর স্বতঃসিদ্ধসমষ্টি গণনসাধ্য। $\Th{PA}$
তত্ত্বটি~$\Th{Q}$-এর তুলনায় অনেক বেশি শক্তিশালী। যেমন, স্বাভাবিক পদ্ধতিতে
আরোহ ব্যবহার করে সহজেই প্রমাণ করা যায় যে যোগ ও গুণ বিনিময়ী। বস্তুত,
বেশির ভাগ সসীমতাবাদী সংখ্যাতাত্ত্বিক ও সংযোজনতাত্ত্বিক যুক্তি
$\Th{PA}$-তে সম্পাদন করা যায়।

$\Th{PA}$ গণনসাধ্যভাবে স্বতঃসিদ্ধায়িত বলে !!{derivability} বিধেয়
$\Prf[\Th{PA}](x,y)$ গণনসাধ্য এবং তাই $\Th{Q}$-তে (এবং সেই কারণে
$\Th{PA}$-তেও) উপস্থাপিত। আগের মতোই $\OPrf[\Th{PA}](x,y)$ দিয়ে
সম্বন্ধটিকে উপস্থাপনকারী সূত্রটি বোঝাব। $\OProv[\Th{PA}](y)$ হলো
$\lexists[x][\OPrf[\Th{PA}](x,y)]$ সূত্রটি; স্বজ্ঞাতভাবে এটি বলে,
``$y$ $\Th{PA}$-এর স্বতঃসিদ্ধগুলি থেকে !!{derivable}।'' $\Th{Q}$-এর
স্বতঃসিদ্ধের চেয়ে সামান্য বেশি যা দরকার, তার কারণ হলো ব্যবহৃত তত্ত্বটি
যেন এই !!{derivability} বিধেয় সম্বন্ধে কয়েকটি মৌলিক তথ্য !!{derive}
করার মতো যথেষ্ট শক্তিশালী হয়। বস্তুত আমাদের নিচের তথ্যগুলি দরকার:
% BN-SRC-244 TARGET-CORRECTION-BEGIN
\par\smallskip\noindent\textnormal{\small\textbf{উৎস-সংশোধন (BN-SRC-244):}
হিমায়িত ইংরেজি সংজ্ঞায় উপস্থাপনকারী বস্তু-ভাষার সূত্র
$\OPrf[\Th{PA}](x,y)$-এর বদলে অধিভাষার সম্বন্ধ
$\Prf[\Th{PA}](x,y)$ বসানো হয়েছে। ঠিক আগের বাক্য ও পরের P1--P3
শর্তের সঙ্গে মিলিয়ে বাংলা পাঠে $\OPrf$ পুনরুদ্ধার করা হয়েছে।}
% BN-SRC-244 TARGET-CORRECTION-END
\begin{enumerate}
\item[P1.] যদি $\Th{PA} \Proves !A$, তবে $\Th{PA} \Proves
  \OProv[\Th{PA}](\gn{!A})$।
\item[P2.] সব !!{formula}s $!A$ ও $!B$-এর জন্য
  \[
  \Th{PA} \Proves \OProv[\Th{PA}](\gn{!A \lif !B}) \lif
  (\OProv[\Th{PA}](\gn{!A}) \lif \OProv[\Th{PA}](\gn{!B})).
  \]
\item[P3.] প্রতিটি !!{formula}~$!A$-এর জন্য
  \[
  \Th{PA} \Proves \OProv[\Th{PA}](\gn{!A})
  \lif \OProv[\Th{PA}](\gn{\OProv[\Th{PA}](\gn{!A})}).
  \]
\end{enumerate}
এই তিনটি ধর্ম সত্য কি না যাচাই করার একমাত্র উপায় হলো
!!{formula}~$\OProv[\Th{PA}](y)$-কে সতর্কভাবে বর্ণনা করা এবং প্রাসঙ্গিক
বিধিবদ্ধ !!{derivation}s বর্ণনা করতে $\Th{PA}$-এর স্বতঃসিদ্ধগুলি ব্যবহার
করা। (১) ও~(২) নং শর্ত সহজ; আসল পরিশ্রম লাগে (৩) নং শর্তে। (কী ধরনের
পরিশ্রম করতে হবে তা ভাবো~\dots) খুঁটিনাটি সম্পূর্ণ করা ক্লান্তিকর এবং
নিরস হবে; তাই এখানে আমরা তোমাকে বিশ্বাস করে নিতে বলব যে $\Th{PA}$-এর
উপরের তিনটি ধর্ম আছে। $\OProv[\Th{PA}](y)$-এর একটি যুক্তিসঙ্গত পছন্দ
নিচেরটিও পূরণ করবে:
\begin{enumerate}
\item[P4.] যদি $\Th{PA} \Proves \OProv[\Th{PA}](\gn{!A})$, তবে
  $\Th{PA} \Proves !A$।
\end{enumerate}
কিন্তু এই তথ্যটি আমাদের লাগবে না।

\begin{digress}
প্রসঙ্গত, এখানে আমরা যেমন খুঁটিনাটি এড়িয়ে যাচ্ছি, গ্যোডেলও তেমনই
করেছিলেন। ১৯৩১ সালের গবেষণাপত্রের শেষে তিনি দ্বিতীয় অসম্পূর্ণতা
উপপাদ্যের প্রমাণের রূপরেখা দেন এবং পরের একটি গবেষণাপত্রে খুঁটিনাটি দেওয়ার
প্রতিশ্রুতি দেন। সে কাজ আর তাঁর হয়ে ওঠেনি; যুক্তিটি যাঁরা বুঝেছিলেন,
তাঁরা সবাই বিশ্বাস করতেন যে সেটি সম্পূর্ণ করা সম্ভব (তাই খুঁটিনাটি পূরণ
করার প্রয়োজন তাঁর হয়নি)।
\end{digress}

\end{document}
